\documentclass{beamer}
\mode<presentation>
\usepackage{amsmath}
\usepackage{amssymb}
\usepackage[utf8]{inputenc}
\usepackage{fontenc}
%\usepackage{advdate}
\usepackage{adjustbox}
\usepackage{subcaption}
\usepackage{enumitem}
\usepackage{multicol}
\usepackage{mathtools}
\usepackage{listings}
\usepackage{url}
\def\UrlBreaks{\do\/\do-}
\usetheme{Boadilla}
\usecolortheme{lily}
\setbeamertemplate{footline}
{
  \leavevmode%
  \hbox{%
  \begin{beamercolorbox}[wd=\paperwidth,ht=2.25ex,dp=1ex,right]{author in head/foot}%
    \insertframenumber{} / \inserttotalframenumber\hspace*{2ex} 
  \end{beamercolorbox}}%
  \vskip0pt%
}
\setbeamertemplate{navigation symbols}{}

\providecommand{\nCr}[2]{\,^{#1}C_{#2}} % nCr
\providecommand{\nPr}[2]{\,^{#1}P_{#2}} % nPr
\providecommand{\mbf}{\mathbf}
\providecommand{\pr}[1]{\ensuremath{\Pr\left(#1\right)}}
\providecommand{\qfunc}[1]{\ensuremath{Q\left(#1\right)}}
\providecommand{\sbrak}[1]{\ensuremath{{}\left[#1\right]}}
\providecommand{\lsbrak}[1]{\ensuremath{{}\left[#1\right.}}
\providecommand{\rsbrak}[1]{\ensuremath{{}\left.#1\right]}}
\providecommand{\brak}[1]{\ensuremath{\left(#1\right)}}
\providecommand{\lbrak}[1]{\ensuremath{\left(#1\right.}}
\providecommand{\rbrak}[1]{\ensuremath{\left.#1\right)}}
\providecommand{\cbrak}[1]{\ensuremath{\left\{#1\right\}}}
\providecommand{\lcbrak}[1]{\ensuremath{\left\{#1\right.}}
\providecommand{\rcbrak}[1]{\ensuremath{\left.#1\right\}}}
\theoremstyle{remark}
\newtheorem{rem}{Remark}
\newcommand{\sgn}{\mathop{\mathrm{sgn}}}
%\providecommand{\abs}[1]{\left\vert#1\right\vert}
\providecommand{\res}[1]{\Res\displaylimits_{#1}} 
\providecommand{\norm}[1]{\lVert#1\rVert}
\providecommand{\mtx}[1]{\mathbf{#1}}
%\providecommand{\mean}[1]{E\left[ #1 \right]}
\providecommand{\fourier}{\overset{\mathcal{F}}{ \rightleftharpoons}}
%\providecommand{\hilbert}{\overset{\mathcal{H}}{ \rightleftharpoons}}
\providecommand{\system}{\overset{\mathcal{H}}{ \longleftrightarrow}}
	%\newcommand{\solution}[2]{\textbf{Solution:}{#1}}
%\newcommand{\solution}{\noindent \textbf{Solution: }}
\providecommand{\dec}[2]{\ensuremath{\overset{#1}{\underset{#2}{\gtrless}}}}
\newcommand{\myvec}[1]{\ensuremath{\begin{pmatrix}#1\end{pmatrix}}}
\let\vec\mathbf

\lstset{
language=C,
frame=single, 
breaklines=true,
columns=fullflexible
}

\numberwithin{equation}{section}

\title{2.10.27}
\author{AI25BTECH11038 -- Tejas Uppala}
\date{}

\begin{document}

% --- Slide Start ---
\begin{frame}[fragile]
\maketitle
\end{frame}
% Question Block
\begin{frame}{Question}
The scalar $\mathbf{A} \cdot ((\mathbf{B} + \mathbf{C}) \times (\mathbf{A} + \mathbf{B} + \mathbf{C}))$ equals\\
    (A) 0\\
    (B) $[\mathbf{A} \mathbf{B} \mathbf{C}]$\\
    (C) $[\mathbf{A} \mathbf{B} \mathbf{C}] \div [\mathbf{B} \mathbf{C} \mathbf{A}]$\\
    (D) None of these
\end{frame}


\begin{frame}{Solution}
The expression is a scalar triple product, which can be written using the box product notation:
\begin{equation}
E = \mathbf{A} \cdot ((\mathbf{B} + \mathbf{C}) \times (\mathbf{A} + \mathbf{B} + \mathbf{C})) = [\mathbf{A} \quad (\mathbf{B} + \mathbf{C}) \quad (\mathbf{A} + \mathbf{B} + \mathbf{C})]    
\end{equation}

The box product is equivalent to the determinant of the matrix formed by the vectors. We can use determinant properties (where rows are the vectors $\mathbf{R}_1, \mathbf{R}_2, \mathbf{R}_3$) to simplify:

\textbf{Row Operation:} Apply the determinant property that allows replacing one row by the sum/difference with other rows.

\textbf{Operation:} Replace the third vector ($\mathbf{R}_3$) by subtracting the first two vectors ($\mathbf{R}_1$ and $\mathbf{R}_2$):

\begin{equation}
    \mathbf{R}_3' = \mathbf{R}_3 - \mathbf{R}_1 - \mathbf{R}_2
\end{equation}
\end{frame}

\begin{frame}
The new third vector is:
\begin{equation}
 \mathbf{R}_3' = (\mathbf{A} + \mathbf{B} + \mathbf{C}) - \mathbf{A} - (\mathbf{B} + \mathbf{C}) = \mathbf{A} + \mathbf{B} + \mathbf{C} - \mathbf{A} - \mathbf{B} - \mathbf{C}  \mathbf{0}   
\end{equation}
The box product becomes:
\begin{equation}
    E = [\mathbf{A} \quad (\mathbf{B} + \mathbf{C}) \quad \mathbf{0}]
\end{equation}

The determinant (or scalar triple product) is \textbf{zero} if one of the vectors is the zero vector ($\mathbf{0}$).
$$E = 0$$
The answer is \textbf{(A) 0}.

\end{frame}
% --- Slide End ---

\end{document}
